% Solutions to the exercises of Chapter 9.
\begin{solutions}
  \item The model is $z_i = x + n_i$, i.e.\ $H = \mathbf 1$ (a column of ones) and $R = \diag(\sigma_i^2)$. The WLS estimate is the BLUE
        (\cref{thm:gauss-markov}):
        \[
          \hat x = \big(\mathbf 1\T R\inv\mathbf 1\big)\inv\mathbf 1\T R\inv z = \frac{\sum_iz_i/\sigma_i^2}{\sum_i1/\sigma_i^2},
          \qquad \Var(\hat x) = \Big(\sum_i\frac{1}{\sigma_i^2}\Big)\inv,
        \]
        the inverse-variance weighted mean. The sample mean $\bar z = \frac1m\sum_iz_i$ is linear and unbiased with variance $\frac{1}{m^2}\sum_i\sigma_i^2$, and by the
        Cauchy--Schwarz inequality $\big(\sum_i\sigma_i^2\big)\big(\sum_i\sigma_i^{-2}\big) \ge m^2$, with equality only when all $\sigma_i$ are equal. So
        $\Var(\bar z) \ge \Var(\hat x)$, strictly unless the variances are equal: then (and only then) the two estimators coincide. One poor sensor can
        ruin a plain average; in the weighted one it simply counts for little.

  \item With $z_i = p_0 + vt_i + n_i$, $H$ has rows $(1, t_i)$ and $\Cov(\hat p_0, \hat v) = (H\T H)\inv$. For $t = -5, \dots, 5$, $\sum t_i = 0$ and $\sum t_i^2 = 110$, so
        \[
          (H\T H)\inv = \begin{bmatrix}11 & 0\\ 0 & 110\end{bmatrix}\inv = \diag\big(0.0909,\ 0.00909\big):
        \]
        $\sigma_{p_0} = 0.30$, $\sigma_v = 0.095$, uncorrelated. For $t = 0, \dots, 10$, $\sum t_i = 55$ and $\sum t_i^2 = 385$, and
        \[
          (H\T H)\inv = \frac{1}{1210}\begin{bmatrix}385 & -55\\ -55 & 11\end{bmatrix} = \begin{bmatrix}0.318 & -0.0455\\ -0.0455 & 0.00909\end{bmatrix}.
        \]
        The velocity variance is unchanged (it depends only on the spread of the times), but $p_0$ is now the value at the \emph{end} of the data,
        an extrapolation, with variance $3.5$ times larger and correlation $-0.85$ with $v$: a line that is too steep must start too low to pass
        through the data. The intercept is best determined at the centroid of the times, where the correlation vanishes.

  \item For $z = Hx + n$ with $n \sim \N(0, R)$, $\log p(z\mid x) = -\tfrac12(z - Hx)\T R\inv(z - Hx) + \text{const}$. Its gradient is
        $\nabla_x\log p = H\T R\inv(z - Hx)$ and its Hessian $\nabla_x^2\log p = -H\T R\inv H$, which does not depend on $z$, so
        $\mathcal I(x) = -\E[\nabla_x^2\log p] = H\T R\inv H$. (The outer-product definition agrees: $\E[H\T R\inv nn\T R\inv H] = H\T R\inv H$.) The matrix bound is
        proved like the scalar one: for an unbiased $\hat x$ and the score $s = \nabla_x\log p$, differentiating $\E[\hat x] = x$ gives $\Cov(\hat x, s) = I$, and
        the covariance matrix of $(\hat x, s)$,
        $\begin{bmatrix}\Cov(\hat x) & I\\ I & \mathcal I\end{bmatrix} \succeq 0$, has the PSD Schur complement $\Cov(\hat x) - \mathcal I\inv \succeq 0$ (Exercise 3.3's argument).
        WLS has covariance $(H\T R\inv H)\inv = \mathcal I\inv$ and attains the bound.

  \item Apply \cref{thm:woodbury} (as in Exercise 3.6) to the covariance of \eqref{eq:bayes-linear}:
        $(P\inv + H\T R\inv H)\inv = P - PH\T(HPH\T + R)\inv HP = P - KHP$ with $K = PH\T(HPH\T + R)\inv$, the covariance of \eqref{eq:lmmse-update}. For the mean,
        multiply out with this expression:
        \[
          (P - KHP)\big(P\inv\mu + H\T R\inv z\big) = \mu - KH\mu + \big(PH\T - KHPH\T\big)R\inv z .
        \]
        Since $K(HPH\T + R) = PH\T$, we have $PH\T - KHPH\T = KR$, so the last term is $Kz$, and the mean is $\mu + K(z - H\mu)$: the two forms agree.

  \item With prior $(\hat x_0, P_0)$, the batch estimate after $k$ measurements has information $P_k\inv = P_0\inv + \sum_{i\le k}H_i\T R_i\inv H_i$ and information
        vector $P_k\inv\hat x_k = P_0\inv\hat x_0 + \sum_{i\le k}H_i\T R_i\inv z_i$. Subtracting the expressions for $k - 1$ gives the recursion in information form
        stated in the chapter. Solving it for $\hat x_k$: $\hat x_k = P_k\big(P_{k-1}\inv\hat x_{k-1} + H_k\T R_k\inv z_k\big)$, and with
        $P_{k-1}\inv = P_k\inv - H_k\T R_k\inv H_k$ this is $\hat x_k = \hat x_{k-1} + P_kH_k\T R_k\inv(z_k - H_k\hat x_{k-1})$. Finally $P_kH_k\T R_k\inv = K_k$ (the
        push-through identity of Exercise 5.6) and $P_k = (I - K_kH_k)P_{k-1}$ by Woodbury. A numerical check with three unknowns, thirty scalar
        measurements of random variances and a weak prior: the RLS estimate and covariance agree with the batch solution after every
        measurement to $3\times10^{-12}$.

  \item The LMMSE estimate uses only second moments: with $\E[x] = \E[z] = 0$, $\Sigma_{xz} = \E[x^4] = \tfrac15$ and $\Sigma_{zz} = \E[x^6] + \sigma^2 = \tfrac17 + \sigma^2$,
        \[
          \hat x_{\mathrm{LMMSE}} = \frac{1/5}{1/7 + \sigma^2}\,z \approx 1.4\,z ,
        \]
        a straight line through a cube-root-shaped relation. The conditional mean $\E[x\mid z]$ follows the cube root, $\approx z^{1/3}$ wherever the noise is
        small compared with $z$. At $x = 0.5$ ($z = 0.125$) the LMMSE estimate is $0.17$, at $x = 0.9$ ($z = 0.73$) it is $1.02$, outside the support of $x$;
        the conditional mean gives $0.50$ and $0.90$. Overall, for $\sigma = 0.01$, the LMMSE root-mean-square error is $0.23$ against $0.075$ for the
        conditional mean (computed numerically). LMMSE does poorly because the dependence of $z$ on $x$ is strongly nonlinear: the best straight line
        fits neither small nor large values. This is the limitation the EKF inherits, mitigated by linearizing locally around the current estimate
        rather than globally.

  \item With $z = Hx + n + \beta$, the estimate $\hat x = Lz$ with $L = (H\T R\inv H)\inv H\T R\inv$ satisfies $LH = I$, so
        \[
          \hat x - x = Ln + L\beta, \qquad \E[\hat x] - x = (H\T R\inv H)\inv H\T R\inv\beta .
        \]
        The noise term has zero mean and covariance $(H\T R\inv H)\inv$, which shrinks as measurements accumulate ($\propto 1/m$ for $m$ comparable measurements).
        The bias term does not shrink: a constant offset in every measurement is averaged, not cancelled. The estimate therefore converges,
        with ever-smaller reported covariance, to the wrong value, and the filter becomes overconfident. Large noise is honest: it is in $R$, so it
        makes the covariance large. An unmodeled bias is invisible to the covariance and, if $\beta$ lies in the range of $H$, invisible in the residuals
        too, because it is indistinguishable from a different $x$. The cure is to model the bias as a state (\cref{ch:processes}), which is exactly
        why IMU biases are part of the filter state.
\end{solutions}
